\documentclass[10pt,a4paper]{amsart}
\usepackage[margin=19mm]{geometry}
\usepackage{amsmath,amssymb,mathtools}
\usepackage[scr=rsfs,cal=euler]{mathalfa}
\usepackage{xfrac}
\usepackage[unicode,hidelinks,bookmarks=false]{hyperref}
\newcommand{\ie}{i.e.\ }
\newcommand{\cf}{cf.\ }
\newcommand{\resp}{respectively\ }
\newcommand{\Subsection}{\subsection}
\providecommand{\nobreakdash}{\nobreak\hbox{-}}
\setlength{\emergencystretch}{2em}
\multlinegap0pt
\usepackage{bm}
\usepackage{bbm}
\usepackage{dsfont}
\usepackage{xspace}
\usepackage{cancel}
\usepackage{mathtools}
\usepackage{amsthm}
\makeatletter
\@ifundefined{lemma}{\newtheorem{lemma}{Lemma}}{}
\@ifundefined{theorem}{\newtheorem{theorem}[lemma]{Theorem}}{}
\makeatother
\newcommand{\speeds}{c_{s,s}}
\title[Accelerated Cops and Robbers]{Accelerated Cops and Robbers}
\author{William B. Kinnersley, Nikolas Townsend}
\date{Source: arXiv:2506.20753v1 (CC BY 4.0)}
\begin{document}
\maketitle

\noindent\emph{Source and licence.} Accelerated Cops and Robbers. Version arXiv:2506.20753v1 (CC BY 4.0), distributed under the Creative Commons Attribution 4.0 International licence, as recorded in the arXiv licence metadata for this exact version and shown on the abstract page https://arxiv.org/abs/2506.20753v1. Full rights notice: licenses/arxiv-2506.20753-CC-BY-4.0.txt.\par\smallskip

\noindent\textit{Editorial scope.} Counted unit: the Theorem whose source label is \texttt{thm:subdivision} in the article (source0.tex lines 147--149, source label \texttt{thm:subdivision}), delivered together with the article's own complete original proof of it (source0.tex lines 150--183, one \texttt{proof} environment of 34 lines). No mathematics is written, repaired or re-derived: statement and proof are the article's own text line for line. Required context retained uncounted: none. Every definition, display and label of those lines is retained unchanged. Omitted from this package and not redistributed: 1--146, 184--1294. Nothing inside a retained excerpt is omitted or altered: every retained body is delivered exactly as the article's lines are written, so no omission receipt of any kind accompanies this package. Citations: the retained excerpts carry no citation command at all, so no bibliography entry of the article is transcribed and no reference is redistributed. Reference adaptation: none was needed and none was made; every reference inside the delivered unit resolves inside it, so no unresolved reference or citation is produced. Printed slips kept literal: no printed word, symbol, hypothesis or number of the counted statement or of its proof was altered. Layout and class: the article's own class is \texttt{amsart} with packages including amsmath, amssymb, amsthm, fullpage, xcolor, ulem, tikz, enumitem; the wrapper is the lane's pinned \texttt{amsart} editorial wrapper at 10pt on A4, which restates the theorem-like environments the retained text needs and adds the article's own macro definitions that the retained text needs (\texttt{\textbackslash speeds}), with extra wrapper packages (bm, bbm, dsfont, xspace, cancel, mathtools). The delivered numbering is the unit's own. Assets: no retained span contains a figure environment, a graphics-inclusion command, a PDF-page inclusion, a table or a TikZ picture; no image file is copied into this package, the package declares no asset dependency and no image file is redistributed with it. Licence: version arXiv:2506.20753v1 of this article is distributed under the Creative Commons Attribution 4.0 International licence, as recorded in the arXiv licence metadata for this exact version and shown on the abstract page https://arxiv.org/abs/2506.20753v1. A full rights notice is delivered at licenses/arxiv-2506.20753-CC-BY-4.0.txt. Characters: every retained excerpt is pure ASCII, so the delivered engine, XeLaTeX with its default Latin Modern text font, reports no missing character.
\begin{theorem}\label{thm:subdivision}
    For every graph $G$ and positive integer $s$, we have $c(G)\le \speeds(G^{(s)})\le c(G)+1$.
\end{theorem}

\begin{proof}
    To show the lower bound, let $k$ speed-$s$ cops play on $G^{(s)}$, where $k<c(G)$.   Since $k < c(G)$, a speed-$1$ robber can evade $k$ speed-$1$ cops on $G$; we show how a speed-$s$ robber can mimic this strategy on $G^{(s)}$ to evade $k$ speed-$s$ cops.
    
    The robber will imagine that he is playing against $k$ cops on $G$, and he will use a winning strategy in the speed-$(1,1)$ game on $G$ to guide his play on $G^{(s)}$.  Throughout the game on $G^{(s)}$, the robber will remain solely on branch vertices.  He will begin the game on the branch vertex in $G^{(s)}$ dictated by his winning strategy on $G$ in the speed-$(1,1)$ game.  Henceforth, in each round he will move a full $s$ steps in order to end his turn on a branch vertex.  Specifically, if the robber's strategy on $G$ prescribes that he should move along edge $xy$ in $G$, then in $G^{(s)}$ the robber will move from $x$ along the subdivided edge (which is a path of length $s$) to $y$.  If a cop occupies a branch vertex $u$ in $G^{(s)}$, then the robber plays as if that cop occupies the vertex $u$ in $G$; if a cop occupies a subdivision vertex in $G^{(s)}$ between branch vertices $u$ and $v$, then the robber plays as if the cop occupies whichever vertex in $\{u,v\}$ the cop passed through most recently.  (If a cop begins the game on a subdivision vertex, then the robber arbitrarily chooses one of the corresponding branch vertices and pretends that the cop occupies that vertex.)
    
    The robber interprets the cop moves in $G^{(s)}$ and translates them to moves in the imagined game as follows.
    \begin{itemize}
        \item 
            If a cop moves from a branch vertex $u$ to a branch vertex $v$ along the subdivided edge between them in $G^{(s)}$, then the robber imagines that the cop moves from $u$ to $v$ in $G$.  %Clearly, this is a legal move in $G$, as the existence of a subdivided edge between $u$ and $v$ in  $G^{(s)}$ requires the existence of edge $uv$ in $G$.
        \item
            If a cop moves from a branch vertex $u$ to a subdivision vertex in $G^{(s)}$, then the robber imagines that the cop stays on the vertex $u$ in $G$.
        \item
            If the cop moves from a subdivision vertex through the branch vertex $u$ to another subdivision vertex, then the robber imagines that the cop moves to or stays on $u$ in $G$.
        \item
            Finally, if a cop moves from a subdivision vertex to a subdivision vertex but does not travel through a branch vertex, then the robber imagines that the cop stays put in $G$.
    \end{itemize}
%    Note that if the cop occupies a subdivision vertex between branch vertices $u$ and $v$ in $G^{(s)}$, then the only branch vertices the cop can reach on her turn are $u$ and $v$.  If the cop occupied $u$, then she would instead be able to reach all branch vertices corresponding to neighbors of $u$ in $G$.  Therefore, a cop occupying a subdivision vertex in $G^{(s)}$ can only prove to be less restrictive for the robber than a cop occupying the branch vertex that she passed through most recently.
    
    On each of the robber's turns, he moves to the branch vertex prescribed by his winning strategy against $k$ cops in $G$.  Since the robber can indefinitely evade $k$ speed-$1$ cops in the imagined game, he can indefinitely evade $k$ cops in the speed-$(s,s)$ game on $G^{(s)}$.
    %\textcolor{blue}{Do I need any more justification here?}

    For the upper bound, we show how $c(G)+1$ cops can capture a robber on $G^{(s)}$.  To do this, the cops will imagine a speed-$(1,1)$ game on $G$ and use a winning strategy for $c(G)$ cops in that game to capture the robber on $G^{(s)}$.

    Let $k=c(G)$, and label the cops $C_1, \dots C_k, C_{k+1}$.  Cops $C_1, \dots, C_k$ will mimic a winning cop strategy on $G$; cop $C_{k+1}$ will (for the moment) move arbitrarily.  For $1 \le i \le k$, if $C_i$ begins on vertex $v$ in the imagined game on $G$, then $C_i$ will begin on the branch vertex $v$ in $G^{(s)}$.  In response, the robber chooses a starting vertex in $G^{(s)}$.  Throughout the game, the cops will ensure that if the robber occupies a branch vertex $v$ in $G^{(s)}$, then he occupies $v$ in the imagined game on $G$; if instead the robber occupies a subdivision vertex between branch vertices $w$ and $x$ in $G^{(s)}$, then he must occupy either $w$ or $x$ in $G$.  To this end, if the robber begins on some branch vertex $w$, then the cops imagine that he occupies $w$ in $G$; if instead he begins on a subdivision vertex between branch vertices $w$ and $x$, then the cops arbitrarily choose either $w$ or $x$ and imagine that the robber occupies that vertex in $G$.

    Henceforth, each of the cops $C_1, \dots, C_k$ uses a winning strategy for $k$ cops in $G$ to guide their play in $G^{(s)}$.  In particular, if the cops' strategy on $G$ calls for $C_i$ to move from vertex $v$ to vertex $w$ in $G$, then $C_i$ moves from $v$ to $w$ in $G^{(s)}$; note that this move must be valid since $v$ and $w$ are distance $s$ apart in $G^{(s)}$.  In this way, cops $C_1, \dots, C_k$ always occupy the same vertices in the game on $G^{(s)}$ as they occupy in the imagined game on $G$.  When the robber moves in $G^{(s)}$, the cops imagine that he moves in the imagined game as follows:
    \begin{itemize}
    \item If the robber does not move to or through a branch vertex in $G^{(s)}$, then the cops imagine that he stays put in $G$.  This is clearly a legal robber move in $G$, and it clearly maintains the condition that if the robber occupies a subdivision vertex in $G^{(s)}$, then he occupies one of the corresponding branch vertices in $G$. 
    
    \item If the robber does move to or through a branch vertex $v$ in $G^{(s)}$, then the cops imagine that he moves to $v$ in $G$.  If the robber began on $v$ in $G^{(s)}$, then he began on $v$ in $G$, so his imagined move in $G$ is legal.  If the robber began on a subdivision vertex in $G^{(s)}$, then he moves to $v$ in $G$.  Since the branch vertices are distance $s$ apart in $G^{(s)}$, the robber can only pass through one branch vertex on a turn, so he must have begun his turn in $G$ on a vertex in $N[v]$; hence his imagined move to $v$ is legal.  Finally, if he began his turn on some other branch vertex $w$ in $G^{(s)}$, then he must have moved from $w$ to $v$ in $G^{(s)}$; consequently the distance in $G^{(s)}$ between $w$ and $v$ must be at most $s$, so we must have $vw \in E(G)$, hence the robber's move from $w$ to $v$ in $G$ is legal.  
    \end{itemize}
    
    Eventually, some cop $C$ captures the robber in the imagined game on $G$. If the robber occupies the corresponding branch vertex $v$ in $G^{(s)}$, then $C$ has captured the robber in $G^{(s)}$ and the cops have won.  Otherwise, the robber must occupy a subdivision vertex between $v$ and some other branch vertex $w$.  For the remainder of the game on $G^{(s)}$, cop $C$ takes $s$ steps toward the robber on each turn.  The remaining $k$ cops imagine that a new game has begun on $G$, in which the robber occupies vertex $w$.  As before, they imitate, on $G^{(s)}$, a winning strategy from $G$.  If the robber ever ends his turn on a branch vertex of $G^{(s)}$, then cop $C$ will capture him.  Moreover, the robber can never move to or through the branch vertex occupied by $C$.  Consequently, cop $C$ will always occupy a branch vertex $x$ in $G^{(s)}$, while in the imagined game on $G$, the robber will occupy some neighbor $y$ of $x$.  Since the cops are following a winning strategy for the imagined game on  $G$, some cop $C'$ eventually captures the robber in that game.  At this point, it must be that the robber occupies a subdivision vertex in $G^{(s)}$ between some branch vertices $x$ and $y$, while cops $C$ and $C'$ occupy both $x$ and $y$.  The robber cannot pass through $x$ or $y$ on his next turn, so on the ensuing cop turn, cop $C$ can capture him.
\end{proof}

\end{document}
